% !PRINTED 49
\weberpage{49}

%<<CONT
\weq{21}{x_1 + x_2 + x_3 = \frac{1}{4}\left(\frac{\sqrt{f(x_2)} - \sqrt{f(x_3)}}{x_2 - x_3}\right)^2,}

und aus (19)

\weq{22}{\sqrt{f(x_1)} = \frac{1}{4}\left(\frac{\sqrt{f(x_2)} - \sqrt{f(x_3)}}{x_2 - x_3}\right)^3 - (x_2+x_3)\frac{\sqrt{f(x_2)} - \sqrt{f(x_3)}}{x_2 - x_3} - \frac{x_3\sqrt{f(x_2)} - x_2\sqrt{f(x_3)}}{x_2 - x_3}.}

\begin{center}
\textbf{\S\,14. \gesperrt{Ursprung der elliptischen Funktionen.}}
\end{center}

Wenn in dem elliptischen Differential erster Gattung in der \gesperrt{Legendreschen} Normalform
\[
\frac{1}{2}\frac{dz}{\sqrt{z(1-z)(1-\varkappa^2 z)}}
\]
die Substitution
\weq{1}{z = \sin^2\varphi}
gemacht und sodann die Integration von $\varphi = 0$ an ausgeführt wird, so entsteht das elliptische Integral erster Gattung
\weq{2}{\int_0^\varphi \frac{d\varphi}{\sqrt{1-\varkappa^2\sin^2\varphi}} = u.}
\gesperrt{Jacobi} nennt die obere Grenze $\varphi$ dieses Integrals seine \gesperrt{Amplitude} und schreibt
\weq{3}{\varphi = \mathrm{am}\,u.}
Die trigonometrischen Funktionen dieses Bogens
\weq{4}{\sin\varphi,\ \cos\varphi,\ \sqrt{1-\varkappa^2\sin^2\varphi} = \Delta\varphi}
sind es nun, die, als Funktionen von $u$ betrachtet, \gesperrt{elliptische Funktionen} genannt und von \gesperrt{Jacobi} mit
\[
\mathrm{sinam}\,u,\quad \mathrm{cosam}\,u,\quad \Delta\,\mathrm{am}\,u,
\]
oder, nach \gesperrt{Gudermann}, kürzer mit
\weq{5}{\mathrm{sn}\,u,\quad \mathrm{cn}\,u,\quad \mathrm{dn}\,u}
bezeichnet werden. Diese Funktionen hängen von den zwei Argumenten $u$ und $\varkappa^2$ ab, und wenn eine genauere Bezeichnung notwendig ist, wird dafür auch
\[
\mathrm{sn}\,(u,\varkappa^2),\quad \mathrm{cn}\,(u,\varkappa^2),\quad \mathrm{dn}\,(u,\varkappa^2)
\]
gesetzt.
